(1) Lorentz方程为 $m\dot v_x=qE+qBv_y$、$m\dot v_y=-qBv_x$，令 $\omega_c=qB/m$、$v_D=E/B$。把速度改写为 $u_x=v_x,u_y=v_y+v_D$，得到 $\dot u_x=\omega_c u_y,\dot u_y=-\omega_c u_x$，可用圆周运动解。

(2) 从静止 $u_x(0)=0,u_y(0)=v_D$ 得 $v_x(t)=v_D\sin\omega_ct$，$v_y(t)=v_D(\cos\omega_ct-1)$。积分得 $\boxed{x(t)=v_D[1-\cos\omega_ct]/\omega_c}$，$\boxed{y(t)=v_D[\sin\omega_ct/\omega_c-t]}$，在平面上形成向负 $y$ 漂移的摆线。

(3) 对任意整周期 $2\pi/\omega_c$，$\langle v_x\rangle=0$、$\langle v_y\rangle=-v_D$，所以 $\mathbf v_{\rm drift}=\mathbf E\times\mathbf B/B^2=-(E/B)\mathbf e_y$。摆线尖点发生在速度同时为零的整数回旋周期，漂移与电荷符号无关。

(4) 动能 $K=m(v_x^2+v_y^2)/2=mv_D^2[1-\cos\omega_ct]$。另一方面 $qEx(t)=qEv_D[1-\cos\omega_ct]/\omega_c=mv_D^2[1-\cos\omega_ct]$，完全相等。因 $\mathbf v\cdot(\mathbf v\times\mathbf B)=0$，磁场不供能；令 $B\to0$ 时展开正弦余弦可恢复 $x\to qEt^2/(2m)$、$y\to0$。
